% LTeX: language=de

\subsection{Grundlagen}

\begin{note}{}
Ein \texthl{\textbf{Zufallsexperiment (ZE)}} ist ein Vorgang / Epxeriment dessen Resultate / Ergebnisse (mindestens 2 verschiedene) zufällig eintreten und das beliebig of wiederholbar ist.\\[-1.0ex]

Alle möglichen \texthl{\textbf{Ergebnisse}} eines ZEs werden in der \texthl{\textbf{Ergebnismenge (Ergebnisraum) $\Omega$}} (\textcolor{myg}{``Omega''}) zusammengefasst: 
\begin{equation*} 
    \Omega = \{ \omega_1; \omega_2; \ldots; \omega_n \} \qquad \qquad \omega_{i}: \quad \text{\textbf{Ergebnisse} des ZEs}
\end{equation*}
Jedes Resultat des ZEs ist dabei genau einem Ergebnis $\omega_i$ zugeordnet.

\begin{Example}{Würfelwurf}{}
    \vspace{-3.0ex}
    \begin{align*}
        \tikzmarknode{aa}{} 
        \overbrace{\Omega_{1}}^{\mathclap{\textcolor{myg}{\text{\footnotesize feinster Ergebnisraum}}}}
        &= \{1; 2; 3; 4; 5; 6\} \tikzmarknode{ab}{} 
        & 
        \overbrace{\abs{\Omega_{1}}}^{\mathclap{\text{\footnotesize\shortstack[l]{\texthl{\textbf{Mächtigkeit}} \textcolor{myg}{von $\Omega_{1}$}\\ \# Elemente}}}}
        &= 6\\[1.0ex]
        \tikzmarknode{ba}{} \Omega_{2} 
        &= \{6; \overbrace{\bar{6}}^{\textcolor{myg}{\text{\footnotesize keine} \ 6}}\} \qquad \tikzmarknode{bb}{}
        & \abs{\Omega_{2}} 
        &= 2 \\[1.0ex]
        \Omega_{3} 
        &= \{\text{gerade}; \text{ungerade}\} 
        & \abs{\Omega_{3}} 
        &= 2 \\[1.0ex]
        \Omega_{4} 
        &= \{\text{gerade}; 1; 3; 5\} 
        & \abs{\Omega_{4}} 
        &= 4 \\[1.0ex]
        \Omega_{5} 
        &= \{\text{gerade}; \text{Primzahl}\} 
        & \textcolor{myr}{\mathrlap{\text{\footnotesize\shortstack[l]{kein Ergebnisraum, da $2$\\ beiden zugeordnet werden kann}}}} &
    \end{align*}
    \begin{tikzpicture}[remember picture, overlay,font=\footnotesize, color=myg]
        \draw[-latex] ([anchor=north west, xshift=0mm, yshift=0mm]aa) coordinate (r)
            to [out=210,in=150,looseness=1] 
            node[midway, anchor=east, xshift=0mm, yshift=0mm] {Vergröberung}
            ([anchor=north west, xshift=0mm, yshift=0mm]ba -| r);
        \draw[-latex] ([anchor=north east, xshift=0mm, yshift=0mm]bb) coordinate (r)
            to [out=30,in=330,looseness=1] 
            node[midway, anchor=west, xshift=0mm, yshift=0mm] {Verfeinerung}
            ([anchor=north east, xshift=0mm, yshift=0mm]ab -| r);
    \end{tikzpicture}
    \vspace{-2.0ex}
\end{Example}
\end{note}